\subsection{Local alternatives need different company histories}

The two asset books expose a constraint on combining explanations. Both fit
the selected PP\&E movements and capital payments. Can they therefore be
inserted interchangeably into one otherwise unchanged history of Apple?
The journals above answer no. Relative to History 1, History 2 retires an
additional 1,000 of gross cost without removing additional accumulated
depreciation, so it recognises an additional 1,000 loss. It acquires an
additional 1,000 of gross cost on credit. Those changes cancel in net PP\&E
and leave cash unchanged, but they add 1,000 to liabilities and subtract
1,000 from profit and equity. Amounts remain millions of dollars, under the
same conditional asset-book assumptions, including exact illustrative
depreciation of 8,200 and no tax effect within that subsystem.

Let $P_i$ denote the vector of statement effects generated by asset History
$i$, and let $R_i$ denote the effects of the rest of the company's history.
For this comparison, order each vector as asset change, liability change,
profit, equity change and cash change. Both full histories must explain the
same reported vector $Y$. If an admissible remainder exists for each branch,
their necessary relation is
\[
P_1+R_1=Y=P_2+R_2,
\qquad
R_2-R_1=-(P_2-P_1)=(0,-1{,}000,+1{,}000,+1{,}000,0).
\]
Thus the rest of History 2 must contribute 1,000 less liability growth and
1,000 more profit and equity growth, with no additional asset or cash change.
Holding the remainder fixed is wrong relative to the common-statement target.
The balance-sheet identity alone fails to detect the error: adding 1,000 to
liabilities and subtracting 1,000 from equity still leaves their sum unchanged.

The reported annual totals make the required differences concrete. Apple's
asset growth is 12,397, liability growth is 17,593, profit is 93,736 and equity
change is $-5{,}196$ \citep[printed pp.~29--32]{apple2024}. Since History 1
contributes liability growth of 1,082, profit of $-8{,}564$ and equity change
of $-8{,}564$, its hypothetical remainder must contribute 16,511, 102,300
and 3,368 respectively. History 2 instead requires 15,511, 103,300 and
4,368. In both cases the remaining asset change must be 19,879. These are
necessary subtraction results. They do not supply the transactions, expense
classifications, taxes or liability assignments that would make a remainder
admissible.

The distinction matters for cash as well. In the isolated asset book,
depreciation and the retirement loss are noncash expenses. Adding both back
to its negative profit gives zero operating cash; acquisitions use 9,447 of
investing cash. Switching branches therefore adds 1,000 to the noncash loss
and 1,000 to the corresponding addback, leaving subsystem cash unchanged.
If that loss were mapped into the annual cash-flow statement's broad Other
adjustment of $-2{,}266$, the remaining contributions to that row would have
to be $-2{,}630$ in History 1 and $-3{,}630$ in History 2
\citep[printed p.~33]{apple2024}. That mapping is itself conditional: the
source does not identify either retirement loss or the composition needed
to offset it. A matching operating-cash total cannot substitute for the
reported adjustment rows.

Nor may a remainder erase the cash-population difference from
Section~\ref{sec:apple-statements}. After the common asset payments, the
required remaining balance-sheet cash change is 9,425, whereas the required
remaining cash-flow change is 8,653. Their 772 difference is inherited from
the disclosed opening populations. Calling it another cash receipt would
manufacture a transaction rather than explain the record.

A proposed completion must now identify the liability class containing the
capital payable, the expense class containing the loss, the tax consequences
and the noncash cash-flow mapping, and satisfy the detailed statement rows
under the same assignments. It must also retain the asset cohorts and payment
terms needed at the next date. The December snapshot does not establish that
either hypothetical annual payable remained outstanding or identify when it
was settled. Failure of the unchanged-remainder construction therefore calls
for a coupled historical model. It rejects that construction without deciding
whether either local asset book has an admissible full-company completion.
